% =====================================================================
\section{Equivalence Classes and Their Limits}
\label{sec:equiv}
% =====================================================================

Two declarations sharing the same content address state the same
mathematical fact (modulo the canonicalization level).  The set of all
declarations with a given address forms an \emph{equivalence class}.

\paragraph{Same hash does not mean interchangeable.}  The content
address captures the \emph{type}, not the proof term, runtime behavior,
or elaboration metadata.  This has important consequences.
\emph{Theorems in direct proofs} are interchangeable, by proof
irrelevance.  \emph{Theorems in tactics} are not: a tactic like
\texttt{simp [Nat.add\_mul]} relies on the \texttt{@[simp]} attribute
attached to that specific \emph{name}; substituting an alias with the
same hash but different attributes would change the tactic's behavior.
\emph{Definitions} are not interchangeable: two definitions with the
same type may compute differently.  \emph{Instances} are not
interchangeable: typeclass instances are registered by name and
participate in instance resolution.

For this reason, resolution is always \emph{by name}.  The content
address is a stable, project-independent key for \emph{discovery} and
\emph{cross-referencing}, but resolution goes through a specific
declaration name that carries the necessary metadata.

\begin{lstlisting}[style=output,caption={Alias discovery via content address}]
Resolved: Nat.add_mul
  Hash:   e1ea8f392f15...  (L1)
  Module: Init.Data.Nat.Basic
  Kind:   theorem
  Type:   forall (n m k : Nat), (n + m) * k = n * k + m * k
  Aliases (1): Nat.right_distrib
\end{lstlisting}

\simone{Move limitations discussion earlier. Currently after practical examples that assume interchangeability. Foreshadow in Section~3.}

\paragraph{Canonicalization and definitional equality.}
The fundamental question underlying the equivalence relation is: which
expressions should receive the same hash?  The ideal answer is ``all
definitionally equal expressions,'' but this is not achievable in
general.  Lean's Calculus of Inductive Constructions (CIC) with proof
irrelevance does not admit unique normal forms under full definitional
equality~\cite{abel2018decidability, coquand1991algorithm}.  The
reduction rules---$\beta$, $\delta$ (at various transparencies),
$\iota$, $\zeta$, and proof irrelevance---do not confluently normalize
all expressions to a unique representative.

Our two-level approach is a pragmatic compromise.  Level~0 captures
syntactic equality modulo universe renaming and metadata---a
conservative but stable equivalence.  Level~1 captures equality modulo
reducible $\delta$-reduction, $\beta$-reduction, $\zeta$-reduction, and
reducible $\iota$-reduction---a strictly coarser equivalence that
collapses common aliases.  The gap between Level~1 and full definitional
equality is large: semireducible and opaque definitions are not
unfolded, so type synonyms at these transparencies create distinct
hashes.  In principle, one could add intermediate levels (semireducible
transparency, instance unfolding), but each level increases computation
cost, reduces stability across Lean versions, and risks divergent
normalization.

\paragraph{Universe polymorphism.}
Lean uses \emph{named} universe polymorphism: declarations are
parameterized over universe level variables with programmer-chosen
names.  Two declarations with identical bodies but different universe
parameter names (e.g., \texttt{.\{u, v\}} vs.\ \texttt{.\{a, b\}})
should hash identically.  Level~0 handles this by renaming universe
parameters to positional indices: the first universe variable
encountered becomes \texttt{u\_0}, the second \texttt{u\_1}, and so on.
The traversal order is deterministic (left-to-right depth-first over
the type expression).


% =====================================================================
